\section{Content at a character and extraction of a primitive class}
\label{sec:extraction}

The boundary calculation in Section~\ref{sec:theta} is useful only if
the theta section itself does not tend to zero at the character being
tested.  We prove this assertion first.  We then lift the sections to
cusp forms and extract a Galois extension.  All precision losses in
this section are bounded independently of the stage.  We require only
a nonzero extension over a characteristic-zero field, not an integral
length formula.

We retain the notation of Section~\ref{sec:theta}.  Write $\chi$
for the universal scalar character over $K$, and $\chi_z$ for
its specialization at a field test $z$.  Its character on $L$
is its restriction, equivalently its norm pullback in the idelic
description.  Write $V'=V\otimes\chi_{D'}$ for the companion plane.

For a field test $z$ and $a\in\{w,\bar w\}$, let
$H^1_{\mathrm{prim},a}(K,V\chi;z)$ mean the degree-one group of the
marked Selmer complex with strict condition at $a$, full condition at
the other place above $p$, and unramified condition at every fixed
place away from $p$.  The moving-place conditions are those of the
deformation; their local complexes are acyclic at the tests in use.
Use the same definition for $V'$.  This defines these groups on the
finite-model field fibers as well as at $z_0$.  A \emph{jump} means
that the indicated group is nonzero.

\begin{proposition}[Theta implication]
\label{extraction:theta-implication}
Let $\mathfrak p$ be a horizontal height-one prime of the deformation
ring, with the auxiliary tuple conditions of Section~\ref{sec:theta}.
If $b_w^0\in\mathfrak p$, then, over $\Frac(R/\mathfrak p)$,
at least one of
\[
 H^1_{\mathrm{prim},a}(K,V\chi),\qquad
 H^1_{\mathrm{prim},a}(K,V'\chi)
\]
is nonzero.  Here $a$ is one fixed side determined by the reciprocity
convention of Section~\ref{sec:theta}.  Replacing $w$ by $\bar w$
replaces $a$ by its conjugate.  Both alternatives use the scalar
character of the determinant problem.
\end{proposition}

We prove the proposition by contradiction.  If both groups vanish
over the divisor field, the open conditions in
Lemma~\ref{theta:test-point} give a point $z_0$ at which both still
vanish, all required local invariants vanish, and the separation
polynomials are nonzero.  Fix this point for the rest of the section.
It is defined over a finite extension $k/\Q_p$; from now on abbreviate
$\chi_{z_0}$ to $\chi$.  The finite characters $\nu_i$ are evaluated
at $z_0$ modulo precisions tending to infinity.

\subsection{A positive coefficient at the chosen point}

Use the auxiliary vector $\varphi$, stabilizer character $d_0$, and
unary nonvanishing degrees $d_{10},d_{20}$ of
Lemma~\ref{theta:auxiliary-data}.  The stabilizer of $e_1\in A$ is
the torus $\mathrm U(Le_2)$.  Write its nonzero period as
\[
 I_{\mathrm{stab}}
 =\int_{[\mathrm U(Le_2)]}\varphi(s)d_0(s)\,ds\ne0.
\]
The one-index unfolding of \cite[Lemma~6.1]{C} concerns the primitive
index $e_1$ and the integral monomial
\[
                  [e,B_1]^{d_1}[e,B_2]^{d_2},
                  \qquad d_1+d_2+1=m.
\]
The brackets are the indicated two-column determinants.  At $Y_e=e_1$
they read the second-row coordinates of $B_1,B_2$.  We now specify
the changes to its finite tests; the integral monomial lattice is
unchanged.

At $p$, put the fixed unary unit test in the second-row $B_1$
coordinate and retain the integer test in $B_2$.  Multiply by the
determinant character in \eqref{eq:source-18}.  At the $j$th moving
prime introduce another copy $v'_j$ of its cyclic parameter.  Factor
the corresponding global character into the slot characters with
parameters $v'_j$ and
$(1+u_j)/(1+v'_j)-1$.  This factorization is global, so it includes
their unramified values at the other places.  Require both heads
$Y_{[e,B_1]}$ and $Y_{[e,B_2]}$ to be integral and invertible, and
insert the respective determinant characters.  All other entries
retain their integral tests.  For the moment these moving tests have
no prescribed target control.

These are bounded integral sums of translates of the original pivot
test.  In fact, write the trailing columns as $Y_{\mathrm{head}}S$,
partition $S$ modulo the moving uniformizer, and subtract
$Y_{\mathrm{head}}S_0$ on a residue class.  Dilating the remaining
trailing coordinates identifies that summand with the pivot test.
The second invertibility condition is a residue condition on $S_0$;
its character factor is the character of its invertible residue
determinant.  All coefficients are integral group elements.  The
scales are powers of the moving residue characteristic, hence
$p$-units.  Although the number of summands may grow, their supremum
norm and the nonarchimedean bound for their sum do not grow.  In
particular no inverse order of a residue unit group occurs.

At an unramified-twist Steinberg place use a line-stabilizer
parahoric on the target and an Iwahori-fixed source vector.  In
split coordinates the test consists of integral columns with the
second coordinate of $Y_e$ zero modulo the uniformizer.  The
integral vector--covector pairs with pairing one and this condition
form a single source-Iwahori orbit: complete the first vector,
whose reduction is in the prescribed line, by a vector in the
kernel of the covector.  The two row-two unary tests are therefore
spherical.  At places kept hyperspecial use the single primitive
norm-one orbit.  At the remaining fixed split places, including
the additive ones, temporarily use the fixed-level tests permitting
the one-index unfolding.

Unfolding now separates a row-one Heisenberg theta function from
\[
 \int_{[\mathrm U(Le_2)]}\varphi(s)\eta^{-1}(s)
            \vartheta_{1,d_1}(sb_1)\vartheta_{2,d_2}(sb_2)\,ds.
\]
Here the determinant character cancels the source character in the
compact integration at $p$.  Thus its increasing conductor does not
shrink that integration domain.  On setting $Y_e=e_1$, its remaining
factor is precisely $\nu_{i,p}$ on the row-two $B_1$ unit coordinate.
The row-one function does not see this factor.  At a moving prime,
the two heads give the two row-two unit cuts, with their separate
slot characters; the primitive orbit volume is unchanged.

Assign the true character to the first unary character, and assign
the two moving slot characters to the corresponding unary
characters.  With the splitting characters of
Lemma~\ref{theta:auxiliary-data} their product is $\eta d_0$.
Averaging in $b_1,b_2$ against the inverse unary characters therefore
gives
\begin{equation}
                   I_{\mathrm{stab}}\,T_1(d_1)T_2(d_2).
\tag{20}\label{eq:source-20}
\end{equation}
Here
\[
 T_j(d)=\int_{[\mathrm U(1)]}\vartheta_{j,d}(s)\xi_j(s)^{-1}\,ds
\]
is the corresponding unary period.  Its assigned scalar parameters
are included in $\xi_j$ and will be displayed when necessary.

This averaging has a fixed denominator.  On local torus units the
weight cancels the characters in the tests, including the deep test
at $p$, so the integrand descends through a fixed compact subgroup.
There are consequently fixed finite torus representatives; by
approximation their $p$-components may be one, and no new moving
component is needed.  For each representative the row-one function
is nonzero at a fixed prime-to-$p$ Jacobi torsion, as in the proof
following the cited unfolding lemma.  The needed Jacobi translations
can be integral at $p$: they preserve $e$ and add multiples of $Y_e$
to the $B$ columns, leaving each determinant twist identically
unchanged.  They preserve the canonical branch and may instead be
included in its Fourier trivialization.  The differential comparison
at these torsions is the fixed prime-to-$p$ isogeny comparison.
Thus a bound for \eqref{eq:source-20} detects the content in the
ordinary Fourier--Jacobi lattice, without degree-dependent divided
powers or conductor-dependent averages.

\subsection{The product character in the unary moment measure}

The extra parameters are needed because the two unary periods need
not be simultaneously nonzero at their respective trivial moving
characters.  We first explain why each period is a nonzero analytic
function even when the true parameter has the prescribed value $t_0$.

All degree variables below lie on the narrow branches fixed in
Lemma~\ref{theta:auxiliary-data}.  A bounded measure on integral
positions gives convergent analytic unit moments on these branches.
For integer degrees tending to infinity, the nonunit positions have
moments tending uniformly to zero.  Hence the degree limits of the
actual $B_2$ integer test are its unit-cut moments.  The same bound
applies after inserting any of the character cuts, by
Lemma~\ref{theta:cut-moments}.  On a closed interior character disk,
the finite group power $(1+t)^{n_i(x)}$ converges uniformly to the
true character: if two exponents agree modulo $p^{a_i}$, their power
ratio tends uniformly to one.  The binomial coefficients are
integral and $a_i\to\infty$.  The degree-power series have uniform
tails after the branch has been narrowed.  We may therefore pass
coefficientwise to limits in the joint degree and character series.

For completeness, the varying coefficient fields cause no
compactness assumption here.  Take bounded arrays in $\C_p$ with
seminorm $\lim_{\mathcal U}|\cdot|_p$, quotient by the arrays of
seminorm zero, and complete the result to a field $\mathcal K$.
An array of positive limiting norm has a bounded inverse on a
filter-large set, which proves the field assertion before
completion.  A nonzero element of $\mathcal K$ records a uniform
upper bound on valuation on a filter-large set of stages.  Whenever
we use such a bound, we restrict to that set and relabel the stages;
all previous limits are preserved.  The moment bounds and uniform series
tails just established justify the analytic limits in $\mathcal K$.
All evaluation parameters chosen below will nevertheless lie in
$\Z_p$ or a fixed finite extension of $k$.

Put the moving slot parameters of $T_1$ equal to zero, and omit
its moving unit cuts for the moment.  There is then a bounded
measure $\mu$ on the true character group $\Gamma\simeq\Z_p$,
with coefficients in a fixed scalar extension, and a nowhere-zero
analytic unit factor $u(d)$, such that its unary period is
$u(d)$ times the evaluation of $\mu$ at
\begin{equation}
                 \widehat{\mathcal C}^{\,d-d_{10}}\nu_t.
\tag{21}\label{eq:source-21}
\end{equation}
The hat denotes the $p$-avatar; the two factors are pulled to $L$
by the norm.  More explicitly write the degree branch as
$d=d_{10}+hn$, $n\in\Z_p$, where the fixed positive integer $h$
kills the finite conductor and torsion factors of
$\widehat{\mathcal C}$.  Put
$\mathcal C_\Gamma=\widehat{\mathcal C}^{\,h}$, a character of $\Gamma$.
The first factor in \eqref{eq:source-21} means $\mathcal C_\Gamma^n$;
we do not require $\widehat{\mathcal C}$ itself to descend to
$\Gamma$.

Here is the measure identity underlying this assertion.  Express
the torus integral using its fixed finite representatives $r$ with
$p$-component one, translate their finite tests to one fixed CM
object, and evaluate all moments in that object's fixed frames.
Concretely, the unary moment construction is linear in the finite
Schwartz test.  For $r$ replace that test by its finite Weil translate,
leaving the CM object and its differential and Igusa frames fixed.
This changes the measure, not the position-power rule or its period
normalization.  The only representative-dependent degree factor
outside that measure is the toric character at $r$.
Let $\mu_r$ be the resulting bounded measure on its two unit
positions $x=(x_\sigma)$.  Absorb the fixed degree $d_{10}$ and all
fixed character factors into $\mu_r$.  The remaining factor of
degree $d-d_{10}$ and true character $t$ is
\[
 \bigl(\widehat{\mathcal C}^{\,d-d_{10}}\nu_t\bigr)(r)^{-1}
 \prod_\sigma x_\sigma^{d-d_{10}}
                   \nu_{t,\sigma}(x_\sigma).
\]
By the weight-character convention and CM reciprocity used in
Section~\ref{sec:theta}, the product on the right is this same
character evaluated at the image $\theta(x)\in\Gamma$ of the
local units.  In this description $r$ in a character argument
denotes the norm-reciprocity image of the torus representative in
$\Gamma$; its finite components outside that quotient were already
absorbed into $\mu_r$.  Thus the displayed factor is
$\bigl(\widehat{\mathcal C}^{\,d-d_{10}}\nu_t\bigr)
(r^{-1}\theta(x))$.  Push each $\mu_r$ forward by
$x\mapsto r^{-1}\theta(x)$ and take their fixed finite weighted
sum to obtain $\mu$.  Its norm is bounded by the fixed moment and
volume bounds.  The coordinate and CM-frame powers that were
removed are common to the moment evaluation, independent of $r$,
and give $u(d)$.  This proves \eqref{eq:source-21} as a measure
identity.  In particular it is not an inference from nonvanishing
at augmentation.  If the local-unit map has image a proper open
subgroup of $\Gamma$, the pushforward construction still gives
exactly the same identity, with the corresponding power map.

Let $\gamma$ be the generator of $\Gamma$ used to define $t$.
The measure transform
$F(Z)=\int_\Gamma(1+Z)^{\ell(g)}\,d\mu(g)$ is a bounded series,
where $g=\gamma^{\ell(g)}$.  At $d=d_{10},t=0$ its value is
nonzero by the auxiliary unary nonvanishing.  Hence $F\ne0$.
For fixed $t=t_0$ its argument, in the branch coordinate $n$, is
\[
                    Z(n)=(1+t_0)\mathcal C_\Gamma(\gamma)^n-1.
\]
The logarithm of $\mathcal C_\Gamma(\gamma)$ on this branch
is nonzero, since its local-unit weight has a nonzero logarithmic
direction.  Therefore $Z(n)$ is nonconstant on every degree disk.
If $F(Z(n))$ vanished there identically, one-variable analytic
uniqueness, applied at an interior accumulation point of the
distinct $Z(n)$, would give $F=0$.  This proves nonvanishing of the
first period as a function of degree at the prescribed $t_0$.
The corresponding assertion for $T_2$, which has no true-character
factor, follows directly from its nonzero value at $d_{20}$.

Restore the moving unit cuts while leaving each period's own
moving character trivial.  Subtracting the uniformizer translate
of an integer unary test expresses each cut by the factors
\[
                  1-\alpha_v\mathcal C(\gamma_v)^{\pm d}.
\]
This is the torus-equivariance calculation of \cite[equation~(6.7)]{C}.
It also includes
the unramified true-character value in $\alpha_v$ for the first
period.  The coefficient $\alpha_v$ is a nonzero $p$-unit,
independent of $d$ for fixed $t$.  The tuple conditions make
$\mathcal C(\gamma_v)$ nontorsion.  Each factor is consequently
a nonzero analytic function of degree, so neither unary period
becomes identically zero.

It follows that, with the original variables fixed at $z_0$, the
two analytic functions
\[
 T_1(a;\mathbf v'),\qquad
 T_2\bigl(-1-a;(1+\mathbf u_0)/(1+\mathbf v')-1\bigr)
\]
are nonzero: for the first use $\mathbf v'=0$; for the second
use $\mathbf v'=\mathbf u_0$.  Their product is nonzero in the
ring of analytic functions on a closed polydisk strictly inside
the open parameter disks and containing both these settings.
Choose $a\in\Z_p$ on its degree branch and $\mathbf v'$ in a
fixed finite extension at which the product is nonzero.  To see
that such choices suffice even over $\mathcal K$, vary one
coordinate at a time over infinite subsets of that fixed field
having interior accumulation points.  Repeated one-variable
uniqueness says that vanishing on their product would force the
series to be zero.

Choose positive integers $d_{1i},d_{2i}$ on the prescribed branches
converging respectively to $a,-1-a$, with both tending to infinity
in the usual order.  At each stage their congruences can be made
arbitrarily accurate and their lower bounds arbitrarily large.
Thus $m_i=d_{1i}+d_{2i}+1$ can simultaneously meet every fixed-stage
weight approximation required by Proposition~\ref{theta:boundary-small}.
The compatibility of the two degree branches was imposed in
Lemma~\ref{theta:auxiliary-data}.  Equation~\eqref{eq:source-20}
then has uniformly bounded valuation at $z_0$.

\begin{lemma}[Content with the local controls]
\label{extraction:content}
There are controlled theta sections at $z_0$, of weights $m_i$
as above and depths $s_i$, whose constant terms tend uniformly to
zero and whose ordinary positive Fourier--Jacobi content has
valuation bounded above by a constant independent of $i$.
\end{lemma}

\begin{proof}
The preceding argument proves the content assertion for the
temporary tests.  Restore the controls at each fixed split place
by the spanning argument of \cite[Lemma~6.4]{C}.  At the auxiliary good
places its finitely many Laurent denominators are nonzero at
$z_0$ by Lemma~\ref{theta:test-point}.  At an additive place the
same local rank filtration works at the present unramified scalar.
A supercuspidal head has no proper Jacquet module.  For a ramified
principal or special head with elliptic determinant trivial on
inertia, both inducing-position characters are ramified.  Neither
can match the unramified lower-rank orbit character.  Thus the
lower-rank coinvariants vanish for every unramified scalar, and
there is no norm linkage between this head and the trailing
unramified binary block.  The open-orbit induction is the
unlinked parabolic product, hence irreducible; a nontrivially
paired pivot generates it.  In a fixed finite-level model choose
a spanning minor nonzero at $z_0$.  Cramer's rule gives the
required expression over a Laurent ring with that single minor
inverted.  This supplies the asserted finite denominators at the
additive places as well.

Cross-multiply these expressions before evaluating the group
algebras.  Their clearing factors have nonzero limits at $z_0$
and bounded inverse valuation.  Prime-to-$p$ target translations
preserve ordinary integrality, since their differential maps are
units and base change of the canonical level is integral.  This
includes the possibly long integral sums at the moving primes:
the nonarchimedean bound for their coefficients is one, and there
is no averaging.  If every controlled section had content tending
to zero, so would every such expression, contradicting
\eqref{eq:source-20}.  We may therefore select a controlled
section with bounded content at each stage.  Only a bounded
number of kinds of fixed local tests is needed.  At every upper
$(2,2)$ control retain the two positive binary operators, the
invertible center expansion and the full last spherical block.
The pivot calculation is unaffected by a determinant twist on
the head.  Finally Proposition~\ref{theta:boundary-small} applies
to all these bounded test variants simultaneously, proving the
constant-term assertion.
\end{proof}

\subsection{Cusp lifting with increasing depth}

We give the depth dependence of the cusp lifting explicitly.
The geometric inputs are \cite[Lemmas~7.1--7.2]{C}: the target has split hyperspecial
PEL datum, multiplicative height three and etale height one at
$p$, and its upper $(2,1,1)$ level marks a rank-two flag and frame
inside the multiplicative rank-three flag.  The last two unit
factors and the upper radical are full.  No splitting of the
etale quotient is marked.  These are exactly our levels $J(s_i)$.

\begin{lemma}[Uniform cusp congruence]
\label{extraction:cusp}
After one fixed reduction of precision, the sections in
Lemma~\ref{extraction:content} give finite integral operator
orders $\mathbb T_i^0$ on classical cusp spaces and maps
\[
 \lambda_i^0:\mathbb T_i^0\longrightarrow
                 \mathcal O_k/(p^{M_i}),\qquad M_i\longrightarrow\infty.
\]
The maps have the tested theta values on the split good and
controlled operators.  Both normalized $p$-operators have value
one.  No upper bound on $s_i$ is required.
\end{lemma}

\begin{proof}
At each depth normalize the good toroidal model in the generic
level, after the finite extension needed for its cyclotomic
frames.  On its canonical ordinary branch the extra level is
finite etale on the Cartier-dual character group of the
multiplicative subgroup.  At a boundary chart this adds a finite
etale cover of the abelian extension-data scheme and leaves the
center Fourier parameter unchanged.  This description holds for
each integer $s_i$; it introduces no division by the degree of
that cover.

A section whose constant terms vanish modulo $p^M$ has positive
Fourier--Jacobi indices.  At such an index the Jacobi line bundle
is ample on the abelian chart.  The differential bundle has its
integral column-degree filtration with constant graded pieces,
so positive cohomology vanishes after tensoring with that ample
line bundle.  This remains true on each finite etale abelian
cover.  Cohomology and base change lift the positive coefficients
modulo $p^M$.  Neatness removes the finite stabilizers of positive
indices, permitting transport of lifts along the free index
orbits without dividing by an orbit order.  Formal functions
then identifies these lifts with the cusp sheaf on the ordinary
open.  This proves the base-change assertion at each $s_i$ with
zero loss of precision, also over ramified finite coefficient
extensions.  Extend by zero on the other ordinary branches.

Multiply by a sufficiently high power of Hasse to extend across
the complement of the ordinary open.  Local lifts of Hasse
differ by a multiple of $p$; sufficiently divisible powers glue
modulo $p^M$ and equal one in multiplicative differential frames.
Serre vanishing after a further ample twist lifts the resulting
global section to characteristic zero.  These exponents may
depend on $i,s_i,M$; the surjectivity statements themselves cost
no precision.  Choose the resulting shift $H_i$ arbitrarily
large and tending to zero on every weight branch in use.  The
classical weights are
\[
             \lambda'_i=(m_i+H_i,m_i+H_i,1+H_i;-1-H_i).
\]

The integral normalized operators remain uniform in depth.
In Serre--Tate coordinates $q_1,q_2,q_3$, each graph for $U_j$
adjoins $j$ roots $u_a^p=q_a$.  Its monomial basis gives
\[
 \Tr\bigl(\mathcal O[u_1,\ldots,u_j]/(u_a^p-q_a)\bigr)
                         \subset p^j\mathcal O.
\]
The transported multiplicative flags and frames depend on the
graph, not on the chosen roots, even when marked modulo
$p^{s_i}$.  They may therefore be taken outside this trace.
The differential scalings have lower bounds
$\lambda'_3+\lambda'_4$ for $j=2$ and $\lambda'_4$ for $j=3$,
exactly as in \cite[Lemma~7.3]{C}.  Thus
\[
 \widetilde U_2=p^{-(2+\lambda'_3+\lambda'_4)}U_2,
 \qquad
 \widetilde U_3=p^{-(3+\lambda'_4)}U_3
\]
preserve the integral lattice.  The graph sum introduces no
denominator.  The Hasse shift changes these exponents by $0,-H_i$;
its determinant factor changes by exactly the same amounts.
Take $H_i$ divisible enough that all remaining unit frame factors
are one modulo the chosen precision.  The two comparison values
are then one.  Prime-to-$p$ operators preserve the lattice and
their comparison values for the same reason.

Use the intersection of the classical cusp space with the
ordinary integral section module as its lattice.  It is a full
bounded lattice and is saturated for divisibility by $p^M$:
its definition as an intersection makes this immediate.  The
commuting integral operators generate an $\mathcal O_k$-algebra
$\mathbb T_i^0$ inside the endomorphisms of that lattice over
a finite model extension of $k$.  This is a finite torsion-free
$\mathcal O_k$-order: the endomorphism lattice is finite over
$\mathcal O_k$, so its submodule generated by the operators is
finite.  The scalar ring of the order is always $\mathcal O_k$;
only its rank and the field representing the lattice may vary
with $i$.  Denote the lifted section by $\mathfrak s_i$.
If $P$ is any integral polynomial relation among these operators,
application to that section gives
\[
               P(\hbox{comparison values})\,\mathfrak s_i
                                      \equiv0\pmod {p^{M_i}}.
\]
Every polynomial has this precision because all its operators
are integral.  Evaluate a positive coefficient of valuation at
most $C$ from Lemma~\ref{extraction:content}.  The scalar
$P(\hbox{comparison values})$ is divisible by $p^{M_i-C}$.
It lies in $k$, so the divisibility descends from the evaluation
field.  This is the length-one version of content cancellation,
performed at $z_0$ itself.  The initial clearing, $C$, and the
neat descent index are fixed; there is no other loss.  Decrease
and relabel $M_i$, also slowing it so that every evaluated finite
group relation holds.  The scalar evaluation factors through
the entire order, giving $\lambda_i^0$.  Its local target selects
the residual local summand of that order, on which both
$\widetilde U_j$ are units.
\end{proof}

\subsection{Local Galois controls and reduction of the order}

The congruence now has to pass to actual cusp eigensystems.
We retain their local subquotients and inertia identities while
removing operator nilpotents, so that these identities survive
in the common limiting matrix algebra.

Apply \cite[Proposition~8.5]{C} at each stage.  Its hypotheses
are the group $\mathrm U(B+\mathbb H)$, discrete holomorphic
infinity of weight $\lambda'_i$ with the two sufficiently large
gaps, the indicated split levels, and hyperspecial level at
places nonsplit over $F$.  These hold here.  It supplies
semisimple four-dimensional representations $\rho_\pi$ with
\[
 \rho_\pi^c\simeq\rho_\pi^\vee\epsilon^{-3},\qquad
 (h_1,h_2,h_3,h_4)
   =(1-m_i-H_i,2-m_i-H_i,2-H_i,1+H_i)
\]
at a distinguished $p$-place.  At split places their
Weil--Deligne parameters have the stated local Langlands
semisimplification and the monodromy upper bound of that
proposition.  They are unramified at nonsplit finite places.
These are characteristic-zero statements at each individual
level.

The positive Jacquet module at $J(s_i)$ allocates unramified
characters to its last two slots.  If their Frobenius roots are
$\alpha_3,\alpha_4$, the normalizations just proved give
\[
 \widetilde U_2=p^{h_3+h_4}/(\alpha_3\alpha_4),\qquad
 \widetilde U_3=p^{h_4}/\alpha_4.
\]
Their unit values give slopes $h_3,h_4$ for the allocated
roots.  The other two slopes sum to $h_1+h_2$.  The sum of the
lowest two slopes is at most this sum and, by weak admissibility,
at least $h_1+h_2$; equality gives the contact at dimension two.
The same argument with the other two roots and $\alpha_3$
gives the contact at dimension three.  The strict Hodge gaps
force strict Newton gaps at these contacts.  On each low-slope
subspace the induced Hodge number is at least the sum of the
corresponding lowest Hodge weights and at most its Newton
number, so equality makes that subspace weakly admissible.
Its canonicity also retains the descent data after a
potentially semistable extension.  Thus
\cite[Lemma~8.6]{C} gives
a stable flag of dimensions $2,3,4$ with last quotient
\[
 X_w=\epsilon^{-h_4}\operatorname{unr}(\widetilde U_3^{-1})
\]
and, at the conjugate place, a first subline
$Y_{\bar w}=X_w^{-c}\epsilon^{-3}$.  The character comparisons
hold on all local elements, uniformly, with limits
$x=\epsilon^{-1}$ and $y=\epsilon^{-2}$.  This argument uses
the two full last unit factors of $J(s_i)$, its positive Jacquet
module, and the weight gaps; none depends on a bound for $s_i$.

We also need binary projections when a Frobenius separation can
vanish at $z_0$.  At a safe auxiliary good place use the full
resultant of \cite[equation~(8.10)]{C}, which excludes ratios by
$Q_v^d$, $|d|\leq4$, between the head and the final binary roots.
At an additive place choose an inertia element $\iota_v$ for
which both comparison head eigenvalues differ from one.  Such
an element exists: elliptic semisimple inertia is nontrivial,
has determinant one, and is finite; in the ramified special
case use an element in the ramifying component.  At a moving
place with nontrivial tested unit character choose a unit on
which it is nontrivial; both head eigenvalues on it are that
scalar.  If the unit character is trivial use the Frobenius
separation retained by Lemma~\ref{theta:test-point}.

For an inertia-tested place let $P_{v,\iota}(Z)$ be the full
degree-four characteristic polynomial.  Its coefficients belong
to the operator order: good split Frobenius coefficients
approximate them simultaneously on the finitely many systems,
and the finite integral operator subalgebra is closed.  Monic
division by $(Z-1)^2$ gives a quadratic $C_{v,\iota}$ in that
order; put $\Delta_v=C_{v,\iota}(1)$.  On each actual system
the remainder is zero, since the allocated unramified final
Jacquet block supplies two inertia eigenvalues one.  The tested
value of $\Delta_v$ is nonzero.  Let $\Delta$ be the product of
these factors and the good-place resultants.

\begin{lemma}[Binary separation at an inertia test]
\label{extraction:binary}
On an actual system with $\Delta_v\ne0$, the generalized
$1$-primary projection $P_v^{\mathrm{bin}}$ of
$\rho_\pi(\iota_v)$ is defined and satisfies
\[
                (\rho_\pi(\iota)-1)P_v^{\mathrm{bin}}=0
                           \qquad(\iota\in I_v).
\]
Together with the Frobenius-separated projections, these identities
give an exponent $C_0$ independent of stage and depth for which
\[
 \Delta^{C_0}\ker(\mathbb T_i^0\longrightarrow\mathbb T_i)=0,
\]
where $\mathbb T_i$ is the image order in characteristic-zero
eigensystems.
\end{lemma}

\begin{proof}
The chosen element has exactly two eigenvalues one, counted
with multiplicity.  Both are already supplied by the final
unramified Jacquet block, so the complementary Weil positions
cannot contain any inertia-trivial constituent.  Thus the
entire inertia-trivial Weil summand is this binary allocation.
An unramified norm twist does not change inertia, so no segment
can cross these allocations.  The segment classification
therefore identifies the allocated final binary representation
with its unramified spherical block, whose predicted monodromy
is zero.  It is necessary to check actual monodromy as well:
the Galois construction in the cited proposition is a direct sum
of shifted copies of the Galois representations of generic
unitary cuspidal factors.  Inside one copy, the selected local
segments have length one; a longer segment would cross the
inertia separation or violate binary sphericity.  Two such
unramified length-one data in a generic unitary factor have real
exponents in an interval of width strictly less than one, so
their quotient cannot be $Q_v$.  Thus monodromy cannot link
them.  It cannot link different copies, which are direct
summands, or link to the head, which is inertia-separated.
The selected subspace consequently has trivial actual inertia.
It has dimension two and is precisely the generalized
$1$-primary space of the chosen element, proving the identity.
The projection can be written explicitly.  Suppress $v$ and
write $P_{\iota_v}(Z)=(Z-1)^2C(Z)$ and $\Delta_v=C(1)$.
The polynomial
\[
 E(Z)=C(Z)\bigl(\Delta_v^{-1}
                  -C'(1)\Delta_v^{-2}(Z-1)\bigr)
\]
is one modulo $(Z-1)^2$ and zero modulo $C(Z)$.
Consequently $P_v^{\mathrm{bin}}=E(\rho_\pi(\iota_v))$.
This gives an algebraic identity with only powers of
$\Delta_v$ in its denominator, suitable for passage to the
limiting corner algebra.  The element $\iota_v$ need not
generate inertia.

For the order, apply the geometric-lemma filtration on each
joint generalized controlled-operator space.  There are at
most $4!$ allocation terms at each controlled place.  Taking
the required compact Levi invariants and a generalized
eigenspace is exact in characteristic zero.  At an
inertia-separated place the two final unramified inducing
positions are forced: no head position can occupy an
unramified slot, because its inertia eigenvalue on
$\iota_v$ differs from one.  At a Frobenius-separated place
the resultants force the allocation.  The full spherical
binary invariants are one-dimensional, so its two operators
are scalars on the surviving grade.  At $p$ the two contacts
force the allocation to the $(2,1,1)$ Levi; both final center
operators are scalars there at every depth.  Permutations
inside the head lie in the Levi Weyl group and do not create
additional grades.  Hence the separated generalized space
has a single grade and no operator nilpotents.

On a component where $\Delta$ has zero eigenvalue, it acts
as zero on every allocation grade and lowers this filtration.
The bounded number of control places therefore gives a fixed
power, for example the product of their bounds $24$, killing
that component.  The good spherical algebra is already
semisimple; hence the kernel of eigensystem detection vanishes
on the separated components and is killed by this power on
the others.  This is an operator identity in characteristic
zero, and the integral order is torsion-free.  It therefore
holds in $\mathbb T_i^0$ itself.
\end{proof}

The comparison values $\lambda_i^0(\Delta)$ have uniformly
bounded valuation, since they tend to the nonzero test values.
Cancellation in Lemma~\ref{extraction:binary} costs at most
$C_0v_p(\lambda_i^0(\Delta))$.  After decreasing precision by
one constant we obtain
\begin{equation}
 \lambda_i:\mathbb T_i\longrightarrow\mathcal O_k/(p^{M_i}),
                      \qquad M_i\longrightarrow\infty.
\label{extraction:reduced-orders}
\end{equation}
At a line-stabilizer parahoric the other local control is
$\rank(\rho_\pi(\iota)-1)\leq1$ and unipotent inertia.
Indeed a single-step flag can separate at most two positions
in a Steinberg block and at most one such block; its predicted
monodromy has rank at most one and square zero.  The Galois
monodromy bound gives the same assertion for actual inertia.
This is also the proof of \cite[Lemma~8.8]{C}, for either orientation
of the line parahoric.

\subsection{The limiting corners and the highest-weight line}

The remaining task is to obtain a nonzero extension of the last
character by the elliptic constituent of the limiting
four-dimensional representation.  We construct the corner
modules encoding these extensions and retain the local identities
just proved.

We use the finite matrix construction of \cite[Section~9,
Lemmas~9.1--9.3]{C} over the field test, namely with
its Artin length equal to one.  We record both its hypotheses
and the resulting objects.  Put $B_i=\mathbb T_i[\rho_i]$, where
$\rho_i$ is the product of all characteristic-zero systems and
stable integral lattices are chosen.  Characteristic coefficients
of all Galois elements lie in $\mathbb T_i$.  The residual
coefficient-prime characteristic polynomials are fixed: reduction
kills every varying group character and every weight-branch
difference.  The ramified supports have bounded cardinality.
After one fixed extension killing the residual semisimplification,
the images are pro-$p$; the global class-field bound for their
number of generators depends on the fixed extension and number
of support places, not their depths.  Cayley--Hamilton together
with the polynomial-identity bound for four-by-four matrices
therefore gives a fixed number of monomials generating $B_i$
over $\mathbb T_i$.

Let $\mathcal A=\prod_{\mathcal U}\mathbb T_i$ and let
$\mathfrak q$ be the kernel of the compact limit
$\mathcal A\to\mathcal O_k$ furnished by
\eqref{extraction:reduced-orders}.  Set
\[
 \mathscr T=(\mathcal A_{\mathfrak q})^h,
 \qquad \lambda:\mathscr T\longrightarrow k,
 \qquad \mathscr B=(\prod_{\mathcal U}B_i)\otimes_{\mathcal A}\mathscr T.
\]
The ring $\mathscr T$ has residue field $k$: its construction
inverts every nonzero constant in $\mathcal O_k$.  Define
\[
 \mathscr D_0=\prod_{\mathcal U}\prod_{\pi\text{ at stage }i}
                         \overline{\Q}_p,
 \qquad
 \mathscr D=\mathscr D_0\otimes_{\mathcal A}\mathscr T.
\]
The embeddings $\mathcal A\hookrightarrow\mathscr D_0$ and
$\prod_{\mathcal U}B_i\hookrightarrow M_4(\mathscr D_0)$
remain injective under the flat localization and henselization.
Thus $\mathscr T\hookrightarrow\mathscr D$ and
$\mathscr B\hookrightarrow M_4(\mathscr D)$ are faithful.
The first ring $\mathscr D_0$ is absolutely flat and reduced;
localization and etale base change preserve these properties,
so $\mathscr D$ is also absolutely flat and reduced.  Its
field evaluations are realized in algebraically closed residue
ultraproducts of the actual systems, with ultrafilters prolonging
$\mathcal U$.  The finite etale choices split in those fields.
These are the ambient realization statements of
\cite[Section~9.1]{C}.
The limiting diagonal blocks are
\begin{equation}
 f=V(-2)\chi^{\pm1},\qquad
 x=\epsilon^{-1},\qquad y=\epsilon^{-2}.
\label{extraction:blocks}
\end{equation}
Here and below $\chi$ is restricted to $G_L$.  These identities
hold on all sequences of group elements: the theta polynomial
comparison holds at split good Frobenius elements,
$\mathcal C^{m_i}$ tends uniformly to one, and at each stage
Frobenius approximation can be taken more accurate than the
chosen precision.  The blocks are absolutely simple and
pairwise distinct; the simplicity of $V|_{G_L}$ was part of
the auxiliary-field choice.
Write $\rho_*$ for the representation in $\mathscr B$ induced
by the arrays $\rho_i$, on the ultraproduct of the stage Galois
groups with their marked local subgroups.

The desired coefficient is $\Hom(y,f)=V\chi^{\pm1}$.
The direct corner from $y$ to $f$ must be nonzero, while paths
through $x$ must not obstruct the triangular quotient.  This
explains the order of the next two lemmas: eliminate the
cyclotomic $y$-to-$x$ extension, then force the $y$-to-$f$
corner to be nonzero by its highest Hodge weight.

The perfect trace pairing on $M_2(k)\times k\times k$ makes
this diagonal representation factor through $\mathscr B$.
Lift the four orthogonal rank-one idempotents by henselian
factorization of an element with four distinct residual
eigenvalues, and lift the inverse matrix units identifying
the two $f$ positions.  For $u,v\in\{f,x,y\}$ let $D_{uv}$
be the finite $\mathscr T$-module of a single-entry corner
from $v$ to $u$.  Diagonal corners are $\mathscr T$, and
\begin{equation}
                 \lambda(D_{uv}D_{vu})=0\quad(u\ne v).
\label{extraction:cycles}
\end{equation}
The corner-minor lemma transports a zero determinant of
stacked matrices to the corresponding corner identity: in
each determinant monomial its directed edges decompose into
one path from the surplus column to the surplus row and
diagonal cycles.  This description also covers two-by-two
minors and is preserved under an algebra map from
$\mathscr B$.

An algebra map from $\mathscr B$ to matrices over $k$ gives
admissible cochains.  Indeed its values on the fixed finite
monomial generators have one common denominator.  Choose
coefficient expressions on the finite quotients of $B_i$,
apply $\lambda_i$, and approximate those fixed matrix values.
This gives continuous matrix functions with increasing uniform
precision and that one denominator.  A relation valid on
all sequences is uniformly valid to increasing precision;
otherwise a violating sequence contradicts it.  The same
argument applies to local splittings.

We need the comparison in the direction from these admissible
cochains to the finite models, not only detection of model
classes by evaluations.  Use the integral contractions
$I_i,Q_i,H_i$ of Lemma~\ref{models:finite-models}, with
$Q_iI_i=1$ and
$1-I_iQ_i=dH_i+H_id$, after the evaluated coefficient base
change modulo the current precision.  Clear the one fixed
denominator in an admissible cocycle $c_i$.  Its cocycle
defect tends uniformly to zero, so $Q_ic_i$ is a model
cocycle to the same increasing precision.  The model has
bounded finite rank over the compact coefficient quotients;
passing to its coefficient limits gives an actual cocycle
over $k$.  The homotopy identity shows that its image under
$I_i$ differs from $c_i$ by the coboundary $dH_ic_i$ and an
error tending to zero.  It therefore represents the original
admissible class.  If the model class were a coboundary, its
image would be an admissible coboundary by the same identity.
Apply the construction to the marked global-local cone,
including its local splitting coordinates.  Its restriction
homotopies preserve every specified local condition.  This
proves the required identification with the field tests.
Thus no bounded local index is an
extra hypothesis of this use of the cited matrix lemmas.

\begin{lemma}[The cyclotomic corner]
\label{extraction:cyclotomic}
The corner modules satisfy $D_{xy}=D_{xf}D_{fy}$.
\end{lemma}

\begin{proof}
Let $Q=D_{xy}/(D_{xf}D_{fy}+\mathfrak mD_{xy})$, where
$\mathfrak m=\ker\lambda$.  A functional $Q\to k$ gives
an extension of $y$ by $x$ with a separate $f$ block; all
discarded products vanish by \eqref{extraction:cycles} and
the definition of $Q$.  If its Galois class splits, the
splitting is preserved by the algebra generated by the
group, hence by the corner idempotents, and the functional
is zero.  We thus obtain an injection into admissible
$H^1(L,k(1))$.

At $\bar w$, stack the matrices
$\rho_*(g)-Y_{\bar w}(g)$.  Their common kernel makes every
four-row determinant zero.  The vanishing of
$\Hom_{G_{L_{\bar w}}}(y,f)$ at $z_0$ and $x\ne y$
supplies three rows with an invertible minor on the competing
columns $f_1,f_2,x$.  Matrix units allow either $f$ coordinate
to be used in either of the two output rows.  Solve these
three equations with $y$ coordinate one.  Each remaining
equation follows from its vanishing four-row corner minor.
This gives a local $y$-eigenline, hence a splitting.  At $w$
apply the transposed argument to the common quotient $X_w$;
the competing minor uses
$\Hom_{G_{L_w}}(f,x)=0$.  These are exactly the local
invariant conditions retained in Lemma~\ref{theta:test-point}.
The extension is therefore strict at every place above $p$.

At a binary place the projection of
Lemma~\ref{extraction:binary}, or its Frobenius version,
is the identity on the $x,y$ extension and makes inertia
trivial there.  At a parahoric place its two-by-two minors
give $(f(\iota)-1)_{ab}u(\iota)=0$ for every entry, where
$u$ is the cyclotomic extension entry.  Some elliptic
inertia difference is nonzero.  Hence $u(\iota)=0$ whenever
that difference is nonzero; multiplying an element in its
kernel by a fixed element outside the kernel and using the
cocycle law gives zero there too.  Elsewhere hyperspecial
unramifiedness applies.  In particular the class is
unramified at the moving places and descends by the
admissible comparison to the fixed field $L$.

Kummer theory identifies the resulting group with $p$-units
tensored with $k$, with zero valuations and logarithms at
all $p$-places.  Zero valuations leave the global units.
The CM biquadratic field $L$ has unit rank one.  A nontrivial
fundamental unit has nonzero $p$-adic logarithm at a split
$p$-place: a unit of $\Q_p$ with zero logarithm is a root
of unity, and the field embedding is injective.  Thus this
one-dimensional space has no element with all local
logarithms zero.  The extension vanishes, so $Q=0$.
Nakayama for the finite module $D_{xy}/D_{xf}D_{fy}$ proves
the assertion.
\end{proof}

\begin{lemma}[Exclusion of the last stable line at arbitrary depth]
\label{extraction:no-line}
The image of $D_{fy}$ is nonzero in every characteristic-zero
field evaluation of the ambient ring.  Consequently
$\operatorname{Fitt}_{0,\mathscr T}(D_{fy})=0$.
\end{lemma}

\begin{proof}
If $D_{fy}$ vanished at such a field, then
Lemma~\ref{extraction:cyclotomic} would make $D_{xy}$ vanish
there as well.  The last idempotent line would be globally
stable.  The competing-column minor used at $\bar w$ is
a unit in $\mathscr T$.  In a matrix with this last line
stable its upper-right column is zero.  The common
$Y_{\bar w}$ kernel, with this competing minor invertible,
must therefore be the last line, so its character is
$Y_{\bar w}$.

These finite matrix relations pass to a filter-large set of
actual systems, as in \cite[proof of Lemma~9.5]{C}.  More explicitly,
stability uses only the bounded algebra generators; the
idempotent, matrix-unit, and competing-minor relations use
finitely many further matrices.  The algebraically closed
residue ultraproduct realizes the henselian choices, so their
finite equalities and inequalities hold on actual systems.
Thus each such system has a global one-dimensional
constituent $\xi$ restricting to $Y_{\bar w}$.

The character $\xi$ is de Rham and hence an algebraic Hecke
character.  Its conjugate weight sums are parallel, by the
relation on the units of the real subfield.  Its weight at
$\bar w$ is $3-h_4$, the isolated lowest weight.  In the
construction of the Galois package it cannot belong to a
factor of Arthur length greater than one: successive
copies would also supply an adjacent integer weight,
whereas the next weight is separated by $h_4-h_3$.
Thus its factor has length one.  The strict unitary generic
bounds on the unramified cuspidal roots place its absolute
Frobenius size strictly between $Q$ and $Q^2$ after the
cohomological shift.  A character of parallel conjugate
weight sum $j\in\Z$ has size $Q^{j/2}$, so $j=3$.

In particular the weight of $\xi$ at $w$ is $h_4$.
Let $F_3\rho_\pi$ be the rank-three local subrepresentation
at the second Hodge--Newton contact.  Its weights are
$h_1,h_2,h_3$, so the $\xi$ line cannot be contained in it:
the de Rham functor is exact and a subrepresentation's
weights occur in this multiset.  Its map to
$\rho_\pi/F_3\rho_\pi=X_w$ is therefore nonzero, and is an
isomorphism of characters.  This identifies the entire
local character of $\xi$ at $w$, including its smooth
factor, without a bound on conductor.  For every local
$g$ we now have the exact identity
\[
 \xi(g)\xi(g^c)
    =X_w(g)Y_{\bar w}(g^c)=\epsilon(g)^{-3}.
\]

Choose one fixed inertia element $g$ with $\epsilon(g)$
nontorsion and include $g,g^c$ in the finite matrix list.
The preceding identity makes
\[
 (e_y\rho_*(g)e_y)(e_y\rho_*(g^c)e_y)
                           -\epsilon(g)^{-3}e_y
\]
zero in the field evaluation.  It is, however, a unit in
$\mathscr T e_y$: its residue under $\lambda$ is
$(\epsilon(g)^{-4}-\epsilon(g)^{-3})e_y\ne0$.
This contradiction proves the nonvanishing of $D_{fy}$.
An element annihilating $D_{fy}$ consequently vanishes in
every ambient field, hence is zero by reducedness and
faithfulness.  The zeroth Fitting ideal of a finite module
is contained in its annihilator by the adjugate identity.
This proves its asserted vanishing.
\end{proof}

\subsection{The extracted class is primitive}

Put $M=D_{fy}\otimes_{\mathscr T,\lambda}k$.  Fitting ideals
commute with base change, so Lemma~\ref{extraction:no-line}
gives $M\ne0$.  A nonzero functional $M\to k$ defines an
extension of $y$ by $f$, keeping $x$ as a separate block.
The only additional path relation is respected because
\[
 D_{fx}D_{xy}=D_{fx}D_{xf}D_{fy}
\]
and the first two factors on the right are a crossing
cycle, killed by \eqref{extraction:cycles}.  The other
discarded products are killed by the diagonal
representation or the same cycle identity.  Thus this is
an algebra representation and yields a nonzero admissible
class with coefficients
\[
                     \Hom(y,f)=V\chi^{\pm1}.
\]
Nonzeroness follows as in Lemma~\ref{extraction:cyclotomic}:
a Galois splitting would be preserved by the full algebra
and its idempotents, forcing the functional to be zero.
The common-kernel argument with rows of types $f,f,x,f$
and the same invertible competing minor makes this
extension split at every conjugate $p$-place.  It has the
full condition at the other half of the $p$-places.

It remains to check unramifiedness at fixed bad places;
this is stronger than the relaxed conclusion used in
\cite[Proposition~9.7]{C}.  At a binary-controlled
place the separated projection is an inertia-fixed
projection whose image maps identically to the $y$
quotient.  Apply it to any lift of that quotient to obtain
an inertia-fixed lift.  The extension class is therefore
unramified there.

At an unramified-twist Steinberg place write the extracted
inertia matrix, after twisting by $y^{-1}$, as
\[
           \begin{pmatrix}F(\iota)&c(\iota)\\0&1\end{pmatrix}.
\]
Here $F=V\chi^{\pm1}$ has its fixed elliptic unipotent
inertia, since the scalar is unramified at this fixed place.
The two-by-two minors on the $f$ rows and the $f,y$ columns
vanish by the parahoric rank bound and the corner-minor
identity.  Consequently $c(\iota)$ belongs to
$\im(F(\iota)-1)$ whenever $F(\iota)\ne1$.
Write $F(\iota)=1+t_p(\iota)N$, where $N\ne0$ has rank
one and $N^2=0$.  Its image $\ell$ is fixed by inertia.
Choose $h$ with $t_p(h)\ne0$.  If $t_p(\iota)=0$, apply
the already proved assertion to $\iota h$ and $h$;
the cocycle law gives $c(\iota)\in\ell$.  Hence all values
of $c$ lie in $\ell$, and $c$ is an additive continuous
homomorphism from inertia to $\ell$.  In residue
characteristic different from $p$, every such homomorphism
factors through the one-dimensional tame pro-$p$ quotient.
Thus $c(\iota)=t_p(\iota)v$ for some $v\in\ell$.
Choose $b$ with $Nb=v$.  Then
$c(\iota)=(F(\iota)-1)b$, an inertia coboundary.
The fixed-local admissible comparison supplies an ordinary
continuous representative, and changing that representative
by a coboundary preserves the argument.  At the other fixed
places hyperspecial unramifiedness applies.  At moving
places retain the allowed full conditions.

Finally descend from $L$ to $K$.  Since $[L:K]=2$ is
invertible in $k$, induction gives
\[
 \operatorname{Ind}_{G_L}^{G_K}(V\chi^{\pm1}|_{G_L})
       \simeq V\chi^{\pm1}\oplus V'\chi^{\pm1}.
\]
This is an isomorphism of the marked local diagrams.  The
strict half of the $p$-places descends from $K$; the full
half does also.  At a fixed non-$p$ place, restriction to
a finite-index inertia subgroup detects unramifiedness in
characteristic zero: restriction followed by corestriction
is multiplication by its nonzero index.  Unramified
inflation contracts any superfluous fixed support.
The same identities pass to our admissible cochains and
finite models.  Applying Shapiro's lemma stagewise and then
passing through the marked comparisons therefore decomposes the
nonzero extracted class into the two primitive one-sided
field tests for $V,V'$; at least one component is nonzero.
If the inverse scalar occurs, complex conjugation identifies
this group with the $\chi$ group on the opposite side, as
fixed in Section~\ref{sec:theta}.  This contradicts the
choice of $z_0$ and proves
Proposition~\ref{extraction:theta-implication}.
